\begin{question}

  \begin{questionpart}
    Prove the following:

    \begin{questionsubpart}
      \begin{equation*}
        \mel{p'}{x}{\alpha} = i\hbar \pdv{}{p'}\braket{p'}{\alpha}
      \end{equation*}
    \end{questionsubpart}

    \begin{questionsubpart}
      \begin{equation*}
        \mel{\beta}{x}{\alpha} = \int dp' \phi_{\beta}^{*}(p')i \hbar \pdv{}{p'} \phi_{\alpha}(p')
      \end{equation*}
    \end{questionsubpart}

    where $\phi_{\alpha}(p') = \braket{p'}{\alpha}$ and
    $\phi_{\beta}(p') = \braket{p'}{\beta}$ are momentum-space wave
    functions.
  \end{questionpart}

  \begin{questionpart}
    What is the physical significance of
    \begin{equation*}
      \exp(\frac{ix\Xi}{\hbar})
    \end{equation*}
  \end{questionpart}

\end{question}

\begin{purpose}
\end{purpose}

\begin{answer}
  \begin{answerpart}{Oh yeah?  Prove it.}

    \begin{answerspart}

      As with most problem in this chapter, we'll start by multiplying
      by 1.

      \begin{equation}
        \begin{split}
          \mel{p'}{x}{\alpha}
          &=(\bra{p'}\hat{x})(\hat{1}\ket{\alpha})\\
          &=\int dx' (\bra{p'}\hat{x})\dyad{x'}(\ket{\alpha})\\
          &=\int dx' \mel{p'}{\hat{x}}{x'}\braket{x'}{\alpha}\\
          &=\int dx' x'\braket{p'}{x'}\braket{x'}{\alpha}\\
        \end{split}
      \end{equation}

      So, here we have two options...we can either postulate the
      existence of some operator $O$ such that our original equation
      holds

      \begin{equation}
        \begin{split}
          \mel{p'}{x}{\alpha}
          &= i\hbar \pdv{}{p'}\braket{p'}{\alpha}\\
          &= \hat{O}\braket{p'}{\alpha}\\
        \end{split}
      \end{equation}

      Substituting that into our earlier equation, we get

      \begin{equation}
        \begin{split}
          \mel{p'}{x}{\alpha}
          &=\int dx' x'\braket{p'}{x'}\braket{x'}{\alpha}\\
          &=\hat{O}\braket{p'}{\alpha}\\
          &=\hat{O}\int dx' (\bra{p'}) \dyad{x'} (\ket{\alpha})\\
          &=\hat{O}\int dx' \braket{p'}{x'}\braket{x'}{\alpha}\\
        \end{split}
      \end{equation}

      So, this is pretty similar to what we got earlier, except that
      we have that pesky operator outside of the integral.  This is
      where we make our second assumption: that it's an operator of
      $p'$ and not $x'$.  If that's the case, then we can migrate the
      operator inside the integral.
      
      \begin{equation}
        \begin{split}
          \mel{p'}{x}{\alpha}
          &=\int dx' x'\braket{p'}{x'}\braket{x'}{\alpha}\\
          &=\int dx' \hat{O}(p')\braket{p'}{x'}\braket{x'}{\alpha}\\
          &=\int dx' x'\braket{p'}{x'}\braket{x'}{\alpha} = \int dx' x'f(x', p')\\
          &=\int dx' \hat{O}(p')\braket{p'}{x'}\braket{x'}{\alpha} = \int dx' \hat{O}(p')f(x', p')\\
        \end{split}
      \end{equation}

      So, clearly, the solution is going to be of the form
      $\exp(kx')$.  Pulling in equation 1.264 because we're going to
      need its complex conjugate in just a moment:

      \begin{equation}
        \tag{1.264}
        \braket{x'}{p'} = \frac{1}{\sqrt{2\pi\hbar}} \exp(\frac{ip'x'}{\hbar})
      \end{equation}

      Separating out the $x'$ and $p'$ portions and substituting in
      equation 1.264 above and we get

      \begin{equation}
        \begin{alignedat}{2}
          & &
          \braket{x'}{\alpha}\hat{O}(p')\braket{p'}{x'}
          &= x'\braket{p'}{x'}\braket{x'}{\alpha}\\
          \implies & &
          \hat{O}(p')\braket{p'}{x'}
          &= x'\braket{p'}{x'}\\
          \implies & &
          \hat{O}(p')\frac{1}{\sqrt{2\pi\hbar}}\exp(\frac{-ip'x'}{\hbar})
          &= x'\frac{1}{\sqrt{2\pi\hbar}}\exp(\frac{-ip'x'}{\hbar})\\
          \implies & &
          \hat{O}(p')\exp(\frac{-ip'x'}{\hbar})
          &= x'\exp(\frac{-ip'x'}{\hbar})\\
        \end{alignedat}
      \end{equation}

      So, we aren't so much deriving this as we are hypothesizing the
      existence and form of it, then computing what it would have to
      be to comply with that hypothesis.  Rewriting the equation
      above, we find

      \begin{equation}
        \begin{alignedat}{2}
          & &
          \hat{O}(p')\exp(\frac{-ip'x'}{\hbar})
          &= x'\exp(\frac{-ip'x'}{\hbar})\\
          \implies & &
          \hat{O}(p')\exp(kp')
          &= x'\exp(kp')\\
          \implies & &
          \hat{O}(p') &= \frac{x'}{k}\pdv{}{p'}\\
          & & &= \cancel{x'}\frac{\hbar}{-i\cancel{x'}}\pdv{}{p'}\\
          & & &= i\hbar\pdv{}{p'}\qed\\
        \end{alignedat}
      \end{equation}

      So, again, we didn't really re-derive it so much as we
      hypothesized the existence of an operator of the form
      $\hat{O}(p)$ acting in the way presented and showed that the
      value is what was presented in the problem.

      As you can probably see by now, the operator we just showed to
      be correct is the analog of the momentum operator we've been
      using the whole time.  You just swap the $p$'s for $x$'s and add
      a negative sign.
    \end{answerspart}

    \begin{answerspart}

      Showing this part is much easier, and also a bit more
      interesting.

      \begin{equation}
        \begin{split}
          \mel{\beta}{x}{\alpha}
          &= \int dp' (\bra{\beta})(\dyad{p'})(x\ket{\alpha})\\
          &= \int dp' \braket{\beta}{p'}\mel{p'}{x}{\alpha}\\
        \end{split}
      \end{equation}

      Using our result from the previous section
      \begin{equation}
        \mel{p'}{x}{\alpha} = i\hbar \pdv{}{p'}\braket{p'}{\alpha}
      \end{equation}

      \begin{equation}
        \begin{split}
          \mel{\beta}{x}{\alpha}
          &= \int dp' \braket{\beta}{p'}\mel{p'}{x}{\alpha}\\
          &= \int dp' \braket{\beta}{p'}\left(i\hbar \pdv{}{p'}\right)\braket{p'}{\alpha}\\
          &= \int dp' \braket{\beta}{p'}\left(i\hbar \pdv{}{p'}\right)\braket{p'}{\alpha}\\
        \end{split}
      \end{equation}

      Pulling in our definitions for $\phi$ from equation 1.265

      \begin{equation}
        \begin{split}
          \phi_{\alpha}(p') &= \braket{p'}{\alpha}\\
          \phi_{\beta}(p') &= \braket{p'}{\beta}\\
        \end{split}
      \end{equation}

      Then substituting those into the equation above

      \begin{equation}
        \begin{split}
          \mel{\beta}{x}{\alpha}
          &= \int dp' \braket{\beta}{p'}\left(i\hbar \pdv{}{p'}\right)\braket{p'}{\alpha}\\          
          &= \int dp' \phi_{\beta}^{*}(p')
          \left(i\hbar \pdv{}{p'}\right)
          \phi_{\alpha}(p')\qed\\
        \end{split}
      \end{equation}

      So, essentially, what we're doing here is looking at the
      \ac{pmf} in momentum space and integrating across it.
    \end{answerspart}
  \end{answerpart}

  \begin{answerpart}
    Before addressing this, it's worth noting the palindrome that
    Sakurai left for us: $ix Xi$...oh what a joker...

    So, we see this equation in a couple of different forms throughout
    our travels and interpreting physically can get hairy.  We're
    going to try to avoid re-creating all of \ac{qft}.  Let's start by
    pulling in the jog operator from section 1.6.4

    \begin{equation}
      \tag{1.218}
      \mathscr{J}(\Delta x' \vu{x}) = \exp(- \frac{ip_x\Delta x'}{N\hbar})
    \end{equation}

    Sign issues aside, that's almost exactly the same as the equation
    we're given:

    \begin{equation}
      \exp(\frac{ix\Xi}{\hbar})
    \end{equation}

    But let's go further.  In addition to being a jog operator, it can
    also be what I'm calling the jerk operator (whether that's in
    honor of Steve Martin or the $3^{rd}$ derivative of position I'll
    leave up to you).

    \begin{equation}
      \mathscr{J}(\Delta p' \vu{p}) = \exp(- \frac{ix\Delta p'}{N\hbar})
    \end{equation}

    We're also applying a phase change in the momentum and position
    spaces.  What does that mean exactly?  Stay tuned for \ac{qft} and
    find out!  Same Feynman time, same Feynman lectures!

    Finally, it also shows us what the wave function would look like
    with a fixed value of $p$ thereby implying an infinite uncertainty
    in position and giving us a plane wave...I figured I'd leave the
    lamest answer for the last one.

  \end{answerpart}

  So...I woooouuullld celebrate the end of the chapter at this point,
  \emph{but}, if you've been reading along then you know that I solved
  this problem \emph{before} solving problem \ref{1.24.theQuestion},
  so look somewhere else for a celebration.

\end{answer}
